\section{Family bounds: finite monodromy and non-effective uniformity}
\label{sec:family-uniform}

\begin{proposition}[R13a: finite Weil monodromy]
Let $f:A\to S$ be a family of polarized abelian fourfolds over a connected
smooth complex variety, with a globally specified $\mathcal O_K$-action of
signature $(2,2)$, where $K=\Q(\omega)$. The rank-two rational local system
of Weil classes becomes constant after a connected finite \'etale cover of
degree at most $6$. After that cover, for any fixed flat rational Weil
class $t$, an integer $\delta\geq1$ with $\delta t$ integral and the number
$T=(t,t)$ can be chosen constant on the entire connected base.
\end{proposition}

\begin{proof}
The integral local system $R^1f_*\Z$ is a torsion-free rank-four module over
the Eisenstein ring $\mathcal O_K$. This ring is a principal ideal domain,
so a fiber is a free $\mathcal O_K$-module. The specified global action
makes monodromy $\mathcal O_K$-linear. Thus, in an integral basis, it lies
in $\operatorname{GL}_4(\mathcal O_K)$. On
\[
 W=\bigwedge\nolimits_K^4 R^1f_*\Q
\]
it acts by its $K$-determinant, which is a unit of $\mathcal O_K$. Since
$\mathcal O_K^\times=\mu_6$, the image has order dividing $6$. The cover
corresponding to its kernel is a finite topological covering. Riemann
existence, SGA~1, Expos\'e~XII, Theorem~5.1, algebraizes it to a finite
\'etale cover of the same degree.

There is also a polarization-based proof of finiteness. The real Weil
plane is primitive of Hodge type $(2,2)$, and its intersection form is
positive definite. The integral lattice
$W\cap R^4f_*\Z(2)$ is preserved by monodromy. The isometry group of a
positive-definite integral lattice is finite: the image of each member
of a fixed basis must lie in a finite set of lattice vectors of a
prescribed norm. This proof does not give the sharper degree-six bound.

After trivializing $W$, a rational vector is a global flat section.
Parallel transport preserves the full integral cohomology lattice, so
an integer clearing its denominator at one point clears it everywhere.
Cup product is flat, so $T$ is constant. It is positive for $t\ne0$.
\end{proof}

\begin{proposition}[R13b: closed threshold loci in a complex family]
Let $f:X\to S$ be smooth and projective, with relatively ample $L$, over
a complex scheme of finite type. Fix a flat rational degree-$2r$ Betti
section $t$. For $C\in\Z_{\geq1}$, the set
\[
 S_C=\{s\in S(\C):c_{L_s}(t_s)\leq C\}
\]
is the set of complex points of a Zariski-closed subset of $S$.
If a single degree bound $D$ for the two integral numerators holds on a
Zariski-dense subset, then a uniform complexity bound holds everywhere.
\end{proposition}

\begin{proof}
Apply the relative Hilbert construction of B01 over $S$ and take the
finite union of the class-correct parameter schemes for $1\leq N\leq C$
and both degrees at most $C$. A relative cycle class is a section of the
pulled-back local system. Equality with $Nt$ holds on a union of connected
components of the parameter scheme; these are algebraic open-and-closed
components. Algebraic and analytic connectedness agree for complex
schemes of finite type, as in SGA~1, XII, Proposition~2.4. The parameter
scheme is projective over $S$, hence its image is closed, and its complex
points are exactly $S_C$.

For bounded numerator degrees without a specified multiplier bound,
apply B06 on each connected component. If $t$ is nonzero there, at most
one multiplier occurs on each connected component of the finite
parameter scheme; if $t=0$, the zero pair suffices. A finite union of
the resulting proper class loci has closed image. A dense subset of
that image forces the image to be all of $S$.
\end{proof}

\begin{theorem}[R13c: pointwise algebraicity implies a family bound]
In the setting of R13b, suppose that $t_s$ is rationally algebraic for
every $s\in S(\C)$. Then there is an integer $C$ such that
\[
 c_{L_s}(t_s)\leq C\qquad\text{for every }s\in S(\C).
\]
The assertion does not provide a numerical value of $C$.
\end{theorem}

\begin{proof}
The closed subsets $S_C$ are nested, and pointwise algebraicity gives
$S(\C)=\bigcup_{C\geq1}S_C$. First suppose $S$ is irreducible and reduced.
If its dimension is zero, the assertion is immediate. Otherwise choose
a small closed Euclidean ball in a complex coordinate chart of its
smooth locus. This ball is a nonempty complete metric space. If every
$S_C$ were a proper Zariski-closed subset of $S$, its intersection with
the ball would be closed and nowhere dense. A proper algebraic subset
has empty analytic interior on the smooth locus: a nonempty analytic
open subset is Zariski dense in an irreducible complex variety. Baire's
theorem forbids the countable union of these nowhere dense sets from
covering the ball. Hence one $S_C$ is all of $S$.

For a general finite-type base, discard nilpotents and apply the argument
to its finitely many irreducible components. Take the maximum of their
finitely many bounds, using that the loci are nested. This yields the
claimed bound on the whole base.
\end{proof}

\begin{remark}[Precise role of the hypotheses]
Finite type over the uncountable field $\C$ is used in the last theorem.
Noetherianity alone does not justify the countable-union argument. A
fixed relatively ample polarization and a fixed flat rational class are
also essential to this statement. It does not give a bound for classes
with arbitrarily growing coefficients, or across infinitely many
polarization types with no common finite-type parameter space.
\end{remark}

\paragraph{Application and the two distinct bounds.}
For a finite-type polarized Eisenstein Weil component, pass to the finite
cover in R13a and fix a rational Weil vector. Any applicable theorem
asserting algebraicity at every complex fiber, such as the precise
Schoen or Markman statements audited elsewhere in this paper, combines
with R13c to give a \emph{non-effective} uniform family bound. This
conclusion is independent of the proposed numerical degree estimate
from the Prym construction.

If a numerical $D_0$ is separately proved on a Zariski-dense subset,
with one effective surface of degree at most $D_0$ and nonzero Weil
projection at every point of that subset, R10--R11 give the explicit
integer
\[
 C=531441\,\delta D_0^4\bigl(1+4\lceil\sqrt T\rceil\bigr).
\]
The constancy of $\delta,T$ and closedness of $S_C$ then extend that
\emph{same numerical bound} to the entire component. Thus the extension
step is proved, while the validity of a proposed numerical $D_0$ remains
a separate source-and-geometry question.

For an arithmetic member, B05 descends a complex representative to its
algebraic field of definition without increasing the complexity, and
B02 specializes that representative at every good place where the
polarization extends. The family bound therefore also bounds these
geometric reductions. This assertion does not turn either a
non-effective family bound or a coarse explicit bound into a
computation of the least complexity at an arbitrary point.

\paragraph{Source check.}
The exact finite-cover statement was read in SGA~1, Expos\'e~XII,
Theorem~5.1, original p.333, recomposed edition p.251, and connectedness
in XII, Proposition~2.4, original p.318, recomposed edition p.243:
\url{https://arxiv.org/pdf/math/0206203}. The full proof of Riemann
existence is retained as background. The closed-locus construction is
proved in B01; the Baire argument and the determinant-monodromy argument
above were independently rederived in this run.
